\section{Scope and Definitions}
This survey focuses on synthetic persona generation for LLM agents. We use
\emph{persona} broadly to mean a representation that conditions an LLM agent as
a user, citizen, customer, patient, student, worker, voter, character,
stakeholder, or other situated actor. A persona may be a structured record, a
natural-language profile, a prompt, a memory state, a retrieval context, or a
hybrid representation. A \emph{persona set} is a collection or distribution of
such representations intended to support some downstream use, such as
simulation, design, evaluation, auditing, or stress testing.

The unit of analysis is the persona-generation procedure: how persona
information is specified, sampled, grounded, enriched, represented, and
validated before or during LLM use. This includes authored design personas,
open-ended model-generated profiles, census- or survey-derived synthetic
populations, trace-grounded user simulations, and hybrid pipelines that combine
data, theory, and LLM rendering. The survey is therefore narrower than work on
LLM agents in general and broader than prompt templates alone.

We distinguish persona generation from the adjacent problem of persona
enactment or adaptation. Training a model to act like a role, user, group, or
individual can improve behavioral fidelity, and this includes supervised
fine-tuning, reinforcement learning, preference optimization, decoding-time
steering, and personalization. However, these methods are not a separate source
of persona information. They can be applied to authored archetypes,
model-generated profiles, population-sampled records, or trace-grounded
histories. We include such work when it evaluates persona following, behavioral
calibration, individual fidelity, safety, or model sensitivity, but we treat it
as an enactment-layer or model-evaluation issue unless it also contributes a
method for constructing the personas themselves.

Several adjacent literatures are therefore in scope only at their boundary with
persona generation: role-playing language agents, LLM-based social simulation,
personalized assistants, synthetic survey respondents, synthetic populations,
and HCI/user-research personas. Work is out of scope when it studies general
agent architectures, general alignment, synthetic data generation, or
personalization without a substantive persona-construction component.
